\documentclass[10pt]{article}

\usepackage[utf8]{inputenc}

\usepackage[T1]{fontenc}

\usepackage{amsmath}

\usepackage{amsfonts}

\usepackage{amssymb}

\usepackage[version=4]{mhchem}

\usepackage{stmaryrd}

\usepackage{graphicx}

\usepackage[export]{adjustbox}

\graphicspath{ {./images/} }

\usepackage{bbold}



\begin{document}

отображении $N$ плотен в $C_{0}$ и $N$ удовлетворяет полному принципу максимума, то существует полугруппа $\left\{P_{t}\right\}$, $t \geqslant 0$, такая, что



\[

N f(x)=\int_{0}^{\infty} P_{t} f(x) d t, \quad f \geqslant 0

\]



(фе л ле ро в ск а я ш олуг р уппа); при этом $P_{t}$ отображает $C_{k}$ в $C_{0}, P_{0}-$ единичное ядро, $\lim _{t \rightarrow 0} P_{t} f=$ $=f, f \in C_{0}$, локально равномерно и $P_{t}(1) \leqslant 1$. Фу0 $f \geqslant 0$ наз. эк с ц е с с и в во й ф у вк ц и е й относительно полугруппы $\left\{P_{t}\right\}$, если всегда $P_{t} f \leqslant f$ и $\lim _{t \rightarrow 0} P_{t} f=$ $=f ;$ если $P_{t} f=f$, то $f$ наз. и н в а р и а н т но о й ц и е й. Соответствующие построения имеют место и для потенциала-меры $\theta N$.



Очерченная теория $\Gamma$. Ханта $(1957-58)$ имеет непосредственный вероятностный смысл. Пусть на $X$ задана нек-рая $\sigma$-алгебра борелевских множеств $\mathfrak{A}$ и вероятностная мера Р. Случайная величина $S=S(x)$ - это $\mathfrak{A}$-измеримое отображение $X$ в пространство состояний $\bar{R}=[-\infty, \infty]$. Семейство случайных величин $\left\{S_{t}\right\}$, $t \geqslant 0,-$ это случайный марковский процесс, для к-рого $S_{t}(x)$ - траектория точки $x \in X$, если для каждого $y$, $-\infty \leqslant y<\infty$, существует вероятностная мера Ру на $\mathfrak{U}$ такая, что



\begin{enumerate}

  \item $\left.\mathrm{P} y\left(\left\{S_{0}=y\right\}\right)=1 ; 2\right) \quad \mathrm{P} y(A), \quad A \in \mathfrak{A}$,

\end{enumerate}



\begin{itemize}

  \item борелевская функция от $y$; 3) вид траектории, проходящей через $y$ в момент $r$, при $t \geqslant r$ не зависит от положения предыдущих ее точек. Применительно к таким марковским процессам полугруппы $\left\{\mathrm{P}_{t}\right\}$ интерпретируются как полугруппы мер

\end{itemize}



\[

\mathrm{P}_{t}(y, B)=\mathrm{P} y\left(\left\{S_{t} \in B\right\}\right) .

\]



Важное значение имеет изучение эксцессивных и инвариантных функций относительно полугрупп $\left\{\mathrm{P}_{t}\right\}$.



C другой стороны, если $X$ есть $\mathfrak{B}$-гармоническое пространство со счетной базой, то на нем всегда можно выбрать ядро потенциалов так, чтобы удовлетворялись условия теоремы Ханта, и эксцессивные функции соответствуюцей полугруппы будут тогда в точности веотрицательными гипергармонич. Функциями. Теорема Ханта обобщается и для нек-рых типов пространств Бауэра (см. [4], [7]).



И другие понятия П. т. а., такие, напр., как выметание, полярные и тонкие множества, также получают вероятностную интерпретацию в рамках общей теории случайных процессов, облегчающую их исследование. C другой стороны, для теории вероятностей оказалась важной теоретико-потенциальная трактовка ряда понятий, таких, напр., как мартингалы, выходящих за рамки марковских процессов.



Jum.: [1] B r e l o t M., Lectures on potential theory, 2 ed., Bombay, 1967; [2] e r o ж e, "L'enseign. math.», 1972, t. 18,

№ 1, p. 1-36; [3] B a e r H., Harmonische Raüme und ihre potentialtheorie, B., 1966 ; [4], Harmonische Raume und inre Co rne a A., potential theory on harmonic spaces, B., 1972; [5] М е й е p П. - А., Вероятность и потенциалы, пер. с англ.,

М., 1973; [6] Х а т Д ж. - А., Марковские процессы и потенМ., 1973; [6] Х а н т Д ж. - А. Марковские процессы и потен-

циалы, пер. с англ., М., 1962; [7] В 1 u m e n t h a l т., G e-



\includegraphics[max width=\textwidth, center]{2023_07_08_534ab5f411755f8ed7c8g-1}

1968.



E. Д. Соломенцев.



ПОТЕНЦИАЛОВ МЕТОД - метод исследования краевых задач для уравнений математич. физики путем сведения их к интегральным уравнениям, основанный на представлении решений этих задач в виде (обобценных) потенциалов.



Пусть в пространстве $\mathbb{R}^{n}, n \geqslant 2$, задано дифференциальное уравнение с частными производными 2-го порядка эллиптич. типа



\[

L u \equiv \sum_{i, j=1}^{n}\left(a_{i j} \frac{\partial^{2} u}{\partial x_{i} \partial x_{j}}\right)+\sum_{i=1}^{n} e_{i} \frac{\partial u}{\partial x_{i}}+c u=f(x)

\]



$=a_{i j}(x), e_{i}=e_{i}(x), c(x) \leqslant 0$ и правой частью $f(x)$, причем $c(x)<-k^{2}<0$ вне нек-рой ограниченной области, содержацей внутри область $D$ класса $C^{1}$. Тогда любое репение $u(x)$ уравнения (1) класса $C^{2}(D \cup S)$ можно представить в виде суммы трех (обобщенных) потенциалов: потенциала объемных масс



\[

\int_{D} E(x, y) \rho(y) d y

\]



потенциала простого слоя



\[

\int_{S} E(x, y) \sigma(y) d s_{y}

\]



и потенциала двойного слоя



\[

\int_{S} Q_{y}[E(x, y)] \mu(y) d s_{y}

\]



где $S=\partial D$ - граница области $D, E(x, y)$ - главное фундаментальное решение оператора $L$, символ $Q_{y}$ обозначает оператор



\[

Q_{y} v=a \frac{\partial v}{\partial N}-b v

\]



действующий в точке $y \in S, N$ - единичный вектор конормали в точке $y \in S$,



\[

\begin{gathered}

a^{2}=\sum_{i=1}^{n}\left(\sum_{j=1}^{n} a_{i j} \cos \left(\nu, y_{j}\right)\right)^{2}, \\

b=\sum_{i=1}^{n} e_{i} \cos \left(\nu, y_{i}\right),

\end{gathered}

\]



$v-$ единичный вектор внешней нормали к $S$ в точке $y \in S$. Плотности потенциалов $\rho(y), \sigma(y)$ и $\mu(y)$ - достаточно гладкие функціи на $D$ или $S$.



Для потенциалов (2)-(4) остаются в силе, с соответствующими изменениями, все дифференциальные и граничные свойства гармонич. потенциалов, описанныө в ст. Потенциала теория для случая, когда $L$ - оператор Лапласа. На основании этих свойств удается свести краевые задачи для эллиптич. уравнений типа (1) к интегральным уравнениям, аналогично тому, как это было ошисапо для задач Дірихле и Неймана для гармонич. функций в ст. lloтенциала теория.



Лum.: [1] М и р а н д а К., Уравнения с частными производными эллиптического типа, пер. с итал., М., $1957 ;[2]$ Б и-

ц а д з е А. В., Краевые задачи для эллиптических уравнений второго порядка, М., 1966; [3] В л а д и м и р о в В. С., Уравнения математической физики, 4 изд., М., 1981 ; [4] К у п р а д3 е В. Д., Методы потенциала в теории упругости, М., 1963 ; [5] М и л н-Т ом с о н Л. М , Теоретическая гидродинамика,

пер. с англ, М , 1964.



ПОТЕНЦИАЛЬНАЯ СЕТЬ, $\quad$ с т ь Е го р ов a,- ортогональная сеть на двумерной поверхности евклидова пространства, к-рую переводит в себя потенциальное двиякение жидкости по этой поверхности. В параметрах П. с. линейный әлемент поверхности имеет вид



\[

d s^{2}=\frac{\partial \Phi}{\partial u} d u^{2}+\frac{\partial \Phi}{\partial v} d v^{2}

\]



где $\Phi=\Phi(u, v)$ - потенциал поля скорости жидкости. Каждая ортогональная полугеодезич. сеть потенциальна. Частным случаем П. с. является Лиувилля сеть. П. с. впервые рассматривал Д. Ф. Егоров (1901).



Лum.: [1] E r о p o в Д. Ф., Работы по дифференциальной геометрии, М., 1970; [2] IÏ у ли и о в с к и й В. И., Классическая дифференциальная геометрия в тензорном изложении,

М., 1963.

В. Базылев.



ПОТЕНЦИАЛЬНОЕ ПОЛЕ, г р а д и е н т н о е п о л е,- векторное поле, образованное градиентами гладкой скалярной функции $f(t)$ нескольких переменных $t=\left(t^{1}, \ldots, t^{n}\right)$, принадлежащих нек-рой области $T$ $n$-мерного пространства. Функция $f(t)$ наз. с к а л я рным потенци а лом (по тенц и а льно й ф ун к ц и е й) этого поля. П. П. вполне интегрируемо в T: П фаффа уравнение (grad $f(t), d t)=0$ имеет в качестве $(n-1)$-мерных интегральных многообразий ли-



c достаточно гладкими коәффициентами $a_{i \jmath}=a_{j i}=$





\end{document}